% =========================
% Explanation: Notation (informal, readable)
% =========================
\subsection*{Explanation of notation (informal)}
We use standard notation for budgeted hyperparameter optimization (HPO) / neural architecture search (NAS).

\begin{itemize}
    \item \textbf{Search space and configurations.}
    \(\Lambda\) denotes the shared search space, i.e., the set of all admissible configurations.
    A single configuration (a hyperparameter assignment or an architecture encoding) is written as \(\lambda \in \Lambda\).
    In practice, \(\lambda\) may be a vector of continuous/integer/categorical hyperparameters (e.g., learning rate, dropout),
    or a structured encoding (e.g., a NAS genotype).

    \item \textbf{Evaluations and objective values.}
    An \emph{evaluation} is one call to the objective function at \(\lambda\), corresponding to one trained model
    (or a benchmark lookup) under the fixed evaluation pipeline.
    In the single-objective case, an evaluation returns a scalar
    \[
        f(\lambda) \in \mathbb{R},
    \]
    such as validation loss or validation error.
    In the multiobjective case, an evaluation returns an \(m\)-dimensional objective vector
    \[
        \mathbf{F}(\lambda) = \big(f_1(\lambda),\ldots,f_m(\lambda)\big) \in \mathbb{R}^m,
    \]
    where \(m\) is the number of objectives (e.g., error, FLOPs, latency, memory).
    Unless stated otherwise, we treat all objectives as \emph{minimization} objectives; any maximization metric is converted to minimization
    (e.g., \(f(\lambda) = -\mathrm{metric}(\lambda)\) or \(f(\lambda)=1-\mathrm{metric}(\lambda)\)).

    \item \textbf{Budget.}
    \(B\) denotes the evaluation budget, i.e., the maximum number of configurations that may be evaluated in one optimization run.
    Thus, each method may perform at most \(B\) evaluations, and comparisons are made under identical \(B\).

    \item \textbf{Stochastic optimization traces and repetitions.}
    Because HPO/NAS methods are typically stochastic, a single run produces a (random) sequence of evaluated configurations
    \[
        \{\lambda_t\}_{t=1}^{B},
    \]
    where \(\lambda_t\) is the configuration evaluated at step \(t\).
    Formally, each optimizer induces a distribution over such length-\(B\) sequences through its internal randomness.
    To estimate performance and variability, we run \(R\) independent repetitions (different random seeds) per method and benchmark instance,
    and aggregate results across these \(R\) runs.
\end{itemize}
